% =========================================================
% 3.8 Boolean Algebra
% File: topics/10_boolean_algebra/topic.tex
% =========================================================

\section{বুলিয়ান অ্যালজেবরা
(\textenglish{Boolean Algebra})}

সাধারণ বীজগণিতে একটি চলক বিভিন্ন সংখ্যামান গ্রহণ করতে পারে।
কিন্তু \textbf{বুলিয়ান অ্যালজেবরায় কোনো চলকের মান কেবল ০ অথবা ১}
হতে পারে। এখানে সাধারণভাবে ০ দ্বারা
\textenglish{False} এবং ১ দ্বারা
\textenglish{True} অবস্থা প্রকাশ করা হয়।

বুলিয়ান অ্যালজেবরার ধারণার প্রবর্তক ছিলেন ইংরেজ গণিতবিদ
\textbf{\textenglish{George Boole}}।
তিনি ১৮৪৭ সালে
\textit{\textenglish{The Mathematical Analysis of Logic}}
এবং ১৮৫৪ সালে
\textit{\textenglish{An Investigation of the Laws of Thought}}
গ্রন্থে যুক্তি ও গণিতের সম্পর্ক নিয়ে আলোচনা করেন।

পরবর্তীতে বুলিয়ান অ্যালজেবরার ধারণা ডিজিটাল ইলেকট্রনিক্স
ও কম্পিউটার সার্কিট বিশ্লেষণের অন্যতম ভিত্তি হিসেবে ব্যবহৃত হতে থাকে।

\par\begingroup

\textbf{সংজ্ঞা:}

যে বীজগাণিতিক পদ্ধতিতে চলকের মান কেবল ০ ও ১ হতে পারে
এবং \textenglish{AND}, \textenglish{OR} ও
\textenglish{NOT}-এর মতো যৌক্তিক ক্রিয়ার মাধ্যমে
বিভিন্ন যৌক্তিক রাশি প্রকাশ ও বিশ্লেষণ করা হয়,
তাকে \textbf{বুলিয়ান অ্যালজেবরা
(\textenglish{Boolean Algebra})} বলে।
\par\endgroup


% =========================================================
% Characteristics
% =========================================================

\subsubsection*{বুলিয়ান অ্যালজেবরার বৈশিষ্ট্য}

\begin{enumerate}[leftmargin=*,itemsep=2pt]
    \item প্রতিটি বুলিয়ান চলকের মান কেবল \textbf{০ অথবা ১} হতে পারে।

    \item প্রধান মৌলিক ক্রিয়া হলো
    \textenglish{OR} \((+)\),
    \textenglish{AND} \((\cdot)\) এবং
    \textenglish{NOT}।

    \item চলকের মানের ক্ষেত্র ০ ও ১-এর মধ্যে সীমাবদ্ধ।

    \item ডিজিটাল সার্কিটের ইনপুট ও আউটপুট বিশ্লেষণে
    এটি ব্যবহৃত হয়।

    \item জটিল বুলিয়ান রাশি ও লজিক সার্কিট সরল করা যায়।

    \item রাশি সরল করলে প্রয়োজনীয়
    \textenglish{Logic Gate}-এর সংখ্যা কমানো সম্ভব হতে পারে।
\end{enumerate}

% =========================================================
% 3.8.1 Boolean Constant and Variable
% =========================================================

\subsection{বুলিয়ান ধ্রুবক ও চলক
(\textenglish{Boolean Constant and Variable})}

\subsubsection*{বুলিয়ান ধ্রুবক (\textenglish{Boolean Constant})}

\par\begingroup

বুলিয়ান অ্যালজেবরায় \textbf{০ ও ১}—এই দুটি নির্দিষ্ট মানকে
\textbf{বুলিয়ান ধ্রুবক
(\textenglish{Boolean Constant})} বলে।
\par\endgroup


অর্থাৎ বুলিয়ান ধ্রুবক দুটি হলো—

\[
0,\qquad 1
\]

\subsubsection*{বুলিয়ান চলক (\textenglish{Boolean Variable})}

\par\begingroup

যে চলক ০ অথবা ১—এই দুই মানের যেকোনো একটি গ্রহণ করতে পারে,
তাকে \textbf{বুলিয়ান চলক
(\textenglish{Boolean Variable})} বলে।
\par\endgroup


উদাহরণ—

\[
A\in\{0,1\}
\]

অর্থাৎ \(A\)-এর মান ০ অথবা ১ হতে পারে।

% =========================================================
% 3.8.2 Boolean vs Ordinary Algebra
% =========================================================

\subsection{বুলিয়ান অ্যালজেবরা ও সাধারণ বীজগণিতের পার্থক্য}

\begin{center}
\begin{tabularx}{\textwidth}{|L{0.22\textwidth}|Y|Y|}
\hline
\rowcolor{ICTLightBlue}
\textbf{বিষয়} &
\textbf{বুলিয়ান অ্যালজেবরা} &
\textbf{সাধারণ বীজগণিত} \\
\hline

চলকের মান &
কেবল ০ বা ১ &
বিভিন্ন সংখ্যামান হতে পারে \\
\hline

প্রধান ক্রিয়া &
\textenglish{AND}, \textenglish{OR}, \textenglish{NOT} &
যোগ, বিয়োগ, গুণ, ভাগ ইত্যাদি \\
\hline

\(1+1\) &
\(1\) &
\(2\) \\
\hline

\(A+A\) &
\(A\) &
\(2A\) \\
\hline

\(A\cdot A\) &
\(A\) &
\(A^2\) \\
\hline

পূরক &
\(A+\overline{A}=1\) &
এ ধরনের নিয়ম প্রযোজ্য নয় \\
\hline

ব্যবহার &
লজিক সার্কিট ও ডিজিটাল সিস্টেম &
সাধারণ গাণিতিক হিসাব \\
\hline
\end{tabularx}
\end{center}

% =========================================================
% 3.8.3 Fundamental Operations
% =========================================================

\subsection{বুলিয়ান অ্যালজেবরার মৌলিক ক্রিয়া}

বুলিয়ান অ্যালজেবরায় তিনটি মৌলিক যৌক্তিক ক্রিয়া হলো—

\begin{enumerate}[leftmargin=*,itemsep=1pt]
    \item \textbf{\textenglish{OR}} — বুলিয়ান যোগ
    \item \textbf{\textenglish{AND}} — বুলিয়ান গুণ
    \item \textbf{\textenglish{NOT}} — বুলিয়ান পূরক
\end{enumerate}

% ---------------------------------------------------------
% OR
% ---------------------------------------------------------

\subsubsection*{বুলিয়ান যোগ বা \textenglish{OR Operation}}

\textenglish{OR} ক্রিয়াকে সাধারণত \(+\) চিহ্ন দ্বারা প্রকাশ করা হয়।

দুটি ইনপুটের \textbf{যেকোনো একটি ১ হলেই ফলাফল ১} হয়।
কেবল উভয় ইনপুট ০ হলে ফলাফল ০ হয়।

\begin{center}
\begin{tabular}{|c|c|c|}
\hline
\rowcolor{ICTLightBlue}
\(A\) & \(B\) & \(A+B\) \\
\hline
0 & 0 & 0 \\
\hline
0 & 1 & 1 \\
\hline
1 & 0 & 1 \\
\hline
1 & 1 & 1 \\
\hline
\end{tabular}
\end{center}

\textbf{\textcolor{ICTOrange}{উদাহরণ:}}

\[
0+0=0,\qquad
0+1=1,\qquad
1+0=1,\qquad
1+1=1
\]

\par\begingroup

\textbf{মনে রাখবে:}

বুলিয়ান যোগ ও বাইনারি গাণিতিক যোগ এক নয়।

\[
\text{\textenglish{Boolean:}}\quad 1+1=1
\]

কিন্তু,

\[
\text{\textenglish{Binary:}}\quad 1+1=10
\]
\par\endgroup


% ---------------------------------------------------------
% AND
% ---------------------------------------------------------

\subsubsection*{বুলিয়ান গুণ বা \textenglish{AND Operation}}

\textenglish{AND} ক্রিয়াকে সাধারণত \(\cdot\) দ্বারা প্রকাশ করা হয়।

\textbf{সব ইনপুট ১ হলেই কেবল ফলাফল ১} হয়।
অন্য যেকোনো ক্ষেত্রে ফলাফল ০।

\begin{center}
\begin{tabular}{|c|c|c|}
\hline
\rowcolor{ICTLightBlue}
\(A\) & \(B\) & \(A\cdot B\) \\
\hline
0 & 0 & 0 \\
\hline
0 & 1 & 0 \\
\hline
1 & 0 & 0 \\
\hline
1 & 1 & 1 \\
\hline
\end{tabular}
\end{center}

\textbf{\textcolor{ICTOrange}{উদাহরণ:}}

\[
0\cdot0=0,\qquad
0\cdot1=0,\qquad
1\cdot0=0,\qquad
1\cdot1=1
\]

% ---------------------------------------------------------
% NOT
% ---------------------------------------------------------

\subsubsection*{বুলিয়ান পূরক বা \textenglish{NOT Operation}}

কোনো বুলিয়ান চলকের মানকে বিপরীত মানে পরিবর্তন করার ক্রিয়াকে
\textbf{\textenglish{NOT Operation}} বা
\textbf{বুলিয়ান পূরক} বলে।

\(A\)-এর পূরক সাধারণত—

\[
\overline{A}
\]

অথবা,

\[
A'
\]

দ্বারা প্রকাশ করা হয়।

অর্থাৎ,

\[
\overline{0}=1
\]

এবং,

\[
\overline{1}=0
\]

\begin{center}
\begin{tabular}{|c|c|}
\hline
\rowcolor{ICTLightBlue}
\(A\) & \(\overline{A}\) \\
\hline
0 & 1 \\
\hline
1 & 0 \\
\hline
\end{tabular}
\end{center}

পূরকের দুটি গুরুত্বপূর্ণ সূত্র হলো—

\[
A+\overline{A}=1
\]

এবং,

\[
A\cdot\overline{A}=0
\]

% =========================================================
% 3.8.4 Precedence
% =========================================================

\subsection{বুলিয়ান ক্রিয়ার অগ্রাধিকার}

একটি বুলিয়ান রাশিতে একাধিক ক্রিয়া থাকলে সাধারণভাবে
নিচের অগ্রাধিকার অনুসরণ করা হয়—

\par\begingroup

\[
\boxed{
\text{\textenglish{Bracket}}
\rightarrow
\text{\textenglish{NOT}}
\rightarrow
\text{\textenglish{AND}}
\rightarrow
\text{\textenglish{OR}}
}
\]
\par\endgroup


\textbf{\textcolor{ICTOrange}{উদাহরণ:}}

\[
A+\overline{B}\cdot C
\]

এখানে প্রথমে \(B\)-এর \textenglish{NOT},
এরপর \textenglish{AND} এবং সবশেষে
\textenglish{OR} ক্রিয়া সম্পন্ন হবে।

% =========================================================
% 3.8.5 Duality
% =========================================================

\subsection{বুলিয়ান দ্বৈতনীতি
(\textenglish{Principle of Duality})}

\par\begingroup

কোনো বৈধ বুলিয়ান রাশিতে \(+\)-এর পরিবর্তে \(\cdot\),
\(\cdot\)-এর পরিবর্তে \(+\), ০-এর পরিবর্তে ১ এবং
১-এর পরিবর্তে ০ বসালে যে নতুন বৈধ বুলিয়ান রাশি পাওয়া যায়,
তাকে মূল রাশির \textbf{দ্বৈত রূপ
(\textenglish{Dual})} বলে।
\par\endgroup


অর্থাৎ—

\[
+\leftrightarrow\cdot
\]

এবং,

\[
0\leftrightarrow1
\]

\textbf{\textcolor{ICTOrange}{উদাহরণ:}}

মূল সূত্র—

\[
A+0=A
\]

এর দ্বৈত রূপ—

\[
A\cdot1=A
\]

আবার,

\[
A+1=1
\]

এর দ্বৈত রূপ—

\[
A\cdot0=0
\]

% =========================================================
% 3.8.6 Boolean Laws
% =========================================================

\subsection{বুলিয়ান সূত্র বা উপপাদ্য}

বুলিয়ান রাশি সরল ও বিশ্লেষণ করার জন্য বিভিন্ন সূত্র ব্যবহৃত হয়।
গুরুত্বপূর্ণ সূত্রগুলো নিচের সারণিতে দেওয়া হলো।

\begin{center}
\small
\begin{tabularx}{\textwidth}{|L{0.27\textwidth}|Y|Y|}
\hline
\rowcolor{ICTLightBlue}
\textbf{সূত্রের নাম} &
\textbf{\textenglish{OR} রূপ} &
\textbf{\textenglish{AND} রূপ} \\
\hline

পরিচিতি সূত্র &
\(A+0=A\) &
\(A\cdot1=A\) \\
\hline

আধিপত্য সূত্র &
\(A+1=1\) &
\(A\cdot0=0\) \\
\hline

অভেদ/পুনরাবৃত্তি সূত্র &
\(A+A=A\) &
\(A\cdot A=A\) \\
\hline

পূরক সূত্র &
\(A+\overline{A}=1\) &
\(A\cdot\overline{A}=0\) \\
\hline

বিনিময় সূত্র &
\(A+B=B+A\) &
\(AB=BA\) \\
\hline

সংযোজন সূত্র &
\(A+(B+C)=(A+B)+C\) &
\(A(BC)=(AB)C\) \\
\hline

বণ্টন সূত্র &
\(A+BC=(A+B)(A+C)\) &
\(A(B+C)=AB+AC\) \\
\hline

শোষণ সূত্র &
\(A+AB=A\) &
\(A(A+B)=A\) \\
\hline

দ্বৈত পূরক সূত্র &
\(\overline{\overline{A}}=A\) &
--- \\
\hline
\end{tabularx}
\end{center}

\subsubsection*{শোষণ সূত্রের উদাহরণ}

প্রমাণ করতে হবে—

\[
A+AB=A
\]

সমাধান—

\[
A+AB=A(1+B)
\]

যেহেতু,

\[
1+B=1
\]

সুতরাং,

\[
A(1+B)=A\cdot1=A
\]

অতএব,

\[
\boxed{A+AB=A}
\]

একইভাবে,

\[
\boxed{A(A+B)=A}
\]

% =========================================================
% 3.8.7 Boolean Addition vs Binary Addition
% =========================================================

\subsection{বুলিয়ান যোগ ও বাইনারি যোগের পার্থক্য}

বুলিয়ান অ্যালজেবরায় \(+\) চিহ্ন দ্বারা
\textenglish{OR Operation} বোঝানো হয়।
কিন্তু বাইনারি গাণিতিক যোগে \(+\) চিহ্ন দ্বারা
গাণিতিক যোগ বোঝানো হয়।

\begin{center}
\begin{tabularx}{\textwidth}{|L{0.25\textwidth}|Y|Y|}
\hline
\rowcolor{ICTLightBlue}
\textbf{বিষয়} &
\textbf{বুলিয়ান যোগ} &
\textbf{বাইনারি যোগ} \\
\hline

\(+\)-এর অর্থ &
\textenglish{OR Operation} &
গাণিতিক যোগ \\
\hline

\(1+1\) &
\(1\) &
\(10\) \\
\hline

\textenglish{Carry} &
হয় না &
হতে পারে \\
\hline

ব্যবহার &
যৌক্তিক ক্রিয়া &
বাইনারি গাণিতিক হিসাব \\
\hline
\end{tabularx}
\end{center}

\par\begingroup

\textbf{ভুল করো না:}

\[
\boxed{
\text{\textenglish{Boolean:}}\quad 1+1=1
}
\]

কিন্তু,

\[
\boxed{
\text{\textenglish{Binary:}}\quad 1+1=10
}
\]
\par\endgroup


% =========================================================
% 3.8.8 Truth Table
% =========================================================

\subsection{সত্যক সারণি
(\textenglish{Truth Table})}

\par\begingroup

কোনো বুলিয়ান রাশি বা লজিক সার্কিটের সম্ভাব্য সকল
ইনপুট সমন্বয়ের বিপরীতে সংশ্লিষ্ট আউটপুটের মান
সারণির মাধ্যমে প্রকাশ করাকে
\textbf{সত্যক সারণি
(\textenglish{Truth Table})} বলে।
\par\endgroup


সত্যক সারণির সাহায্যে—

\begin{enumerate}[leftmargin=*,itemsep=2pt]
    \item লজিক সার্কিটের আউটপুট নির্ণয় করা যায়।
    \item বুলিয়ান সূত্র যাচাই করা যায়।
    \item দুটি বুলিয়ান রাশির সমতুল্যতা যাচাই করা যায়।
    \item সার্কিটের কার্যপ্রণালি বোঝা যায়।
    \item সরলীকৃত রাশির ফল যাচাই করা যায়।
\end{enumerate}

\subsubsection*{সত্যক সারণিতে সারির সংখ্যা}

যদি কোনো বুলিয়ান রাশিতে \(n\)টি ইনপুট চলক থাকে,
তবে সম্ভাব্য ইনপুট সমন্বয়ের সংখ্যা—

\[
\boxed{2^n}
\]

\textbf{\textcolor{ICTOrange}{উদাহরণ:}}

দুটি চলক \(A,B\) থাকলে—

\[
2^2=4
\]

টি সম্ভাব্য ইনপুট সমন্বয় হবে।

\begin{center}
\begin{tabular}{|c|c|}
\hline
\rowcolor{ICTLightBlue}
\(A\) & \(B\) \\
\hline
0 & 0 \\
\hline
0 & 1 \\
\hline
1 & 0 \\
\hline
1 & 1 \\
\hline
\end{tabular}
\end{center}

আবার তিনটি চলক থাকলে—

\[
2^3=8
\]

টি সম্ভাব্য ইনপুট সমন্বয় হবে।

% =========================================================
% 3.8.9 De Morgan's Theorem
% =========================================================

\subsection{ডি-মরগ্যানের উপপাদ্য
(\textenglish{De Morgan's Theorem})}

গণিতবিদ \textbf{\textenglish{Augustus De Morgan}}-এর নামে পরিচিত
দুটি গুরুত্বপূর্ণ উপপাদ্য বুলিয়ান রাশি ও
লজিক সার্কিট সরলীকরণে ব্যাপকভাবে ব্যবহৃত হয়।

ডি-মরগ্যানের উপপাদ্য দুটি হলো—

\[
\boxed{
\overline{A+B}
=
\overline{A}\cdot\overline{B}
}
\]

এবং,

\[
\boxed{
\overline{AB}
=
\overline{A}+\overline{B}
}
\]

% =========================================================
% First Theorem
% =========================================================

\subsubsection*{প্রথম উপপাদ্য}

\par\begingroup

দুই বা ততোধিক চলকের
\textenglish{OR}-এর পূরক তাদের পৃথক পূরকের
\textenglish{AND}-এর সমান।
\par\endgroup


দুটি চলকের ক্ষেত্রে—

\[
\boxed{
\overline{A+B}
=
\overline{A}\cdot\overline{B}
}
\]

\subsubsection*{সত্যক সারণির সাহায্যে প্রথম উপপাদ্য যাচাই}

\begin{center}
\scriptsize
\resizebox{\textwidth}{!}{%
\begin{tabular}{|c|c|c|c|c|c|c|}
\hline
\rowcolor{ICTLightBlue}
\(A\) &
\(B\) &
\(A+B\) &
\(\overline{A+B}\) &
\(\overline{A}\) &
\(\overline{B}\) &
\(\overline{A}\cdot\overline{B}\) \\
\hline
0 & 0 & 0 & 1 & 1 & 1 & 1 \\
\hline
0 & 1 & 1 & 0 & 1 & 0 & 0 \\
\hline
1 & 0 & 1 & 0 & 0 & 1 & 0 \\
\hline
1 & 1 & 1 & 0 & 0 & 0 & 0 \\
\hline
\end{tabular}%
}
\end{center}

সারণিতে প্রতিটি ইনপুট সমন্বয়ের জন্য
\(\overline{A+B}\) এবং
\(\overline{A}\cdot\overline{B}\)-এর মান সমান।

অতএব,

\[
\boxed{
\overline{A+B}
=
\overline{A}\cdot\overline{B}
}
\]

যাচাই হলো।

% =========================================================
% Second Theorem
% =========================================================

\subsubsection*{দ্বিতীয় উপপাদ্য}

\par\begingroup

দুই বা ততোধিক চলকের
\textenglish{AND}-এর পূরক তাদের পৃথক পূরকের
\textenglish{OR}-এর সমান।
\par\endgroup


দুটি চলকের ক্ষেত্রে—

\[
\boxed{
\overline{AB}
=
\overline{A}+\overline{B}
}
\]

\subsubsection*{সত্যক সারণির সাহায্যে দ্বিতীয় উপপাদ্য যাচাই}

\begin{center}
\scriptsize
\resizebox{\textwidth}{!}{%
\begin{tabular}{|c|c|c|c|c|c|c|}
\hline
\rowcolor{ICTLightBlue}
\(A\) &
\(B\) &
\(AB\) &
\(\overline{AB}\) &
\(\overline{A}\) &
\(\overline{B}\) &
\(\overline{A}+\overline{B}\) \\
\hline
0 & 0 & 0 & 1 & 1 & 1 & 1 \\
\hline
0 & 1 & 0 & 1 & 1 & 0 & 1 \\
\hline
1 & 0 & 0 & 1 & 0 & 1 & 1 \\
\hline
1 & 1 & 1 & 0 & 0 & 0 & 0 \\
\hline
\end{tabular}%
}
\end{center}

সারণিতে প্রতিটি ইনপুট সমন্বয়ের জন্য
\(\overline{AB}\) এবং
\(\overline{A}+\overline{B}\)-এর মান সমান।

অতএব,

\[
\boxed{
\overline{AB}
=
\overline{A}+\overline{B}
}
\]

যাচাই হলো।

% =========================================================
% Three-variable De Morgan
% =========================================================

\subsubsection*{তিনটি চলকের ক্ষেত্রে ডি-মরগ্যানের উপপাদ্য}

ডি-মরগ্যানের উপপাদ্য শুধু দুটি চলকের জন্য নয়;
তিন বা ততোধিক চলকের ক্ষেত্রেও একই নিয়ম প্রযোজ্য।

তিনটি চলক \(A,B,C\) হলে প্রথম উপপাদ্য—

\[
\boxed{
\overline{A+B+C}
=
\overline{A}\cdot
\overline{B}\cdot
\overline{C}
}
\]

এবং দ্বিতীয় উপপাদ্য—

\[
\boxed{
\overline{ABC}
=
\overline{A}+
\overline{B}+
\overline{C}
}
\]

তিনটি ইনপুট চলক থাকায় সম্ভাব্য ইনপুট সমন্বয়ের সংখ্যা—

\[
2^3=8
\]

টি।

% =========================================================
% Three-variable First Theorem Truth Table
% =========================================================

\subsubsection*{তিন চলকের প্রথম উপপাদ্য যাচাই}

প্রমাণ করতে হবে—

\[
\overline{A+B+C}
=
\overline{A}\cdot
\overline{B}\cdot
\overline{C}
\]

\begin{center}
\scriptsize
\resizebox{\textwidth}{!}{%
\begin{tabular}{|c|c|c|c|c|c|}
\hline
\rowcolor{ICTLightBlue}
\(A\) &
\(B\) &
\(C\) &
\(A+B+C\) &
\(\overline{A+B+C}\) &
\(\overline{A}\cdot\overline{B}\cdot\overline{C}\) \\
\hline
0 & 0 & 0 & 0 & 1 & 1 \\
\hline
0 & 0 & 1 & 1 & 0 & 0 \\
\hline
0 & 1 & 0 & 1 & 0 & 0 \\
\hline
0 & 1 & 1 & 1 & 0 & 0 \\
\hline
1 & 0 & 0 & 1 & 0 & 0 \\
\hline
1 & 0 & 1 & 1 & 0 & 0 \\
\hline
1 & 1 & 0 & 1 & 0 & 0 \\
\hline
1 & 1 & 1 & 1 & 0 & 0 \\
\hline
\end{tabular}%
}
\end{center}

শেষের দুটি কলামের মান প্রতিটি সারিতে সমান।

সুতরাং,

\[
\boxed{
\overline{A+B+C}
=
\overline{A}\cdot
\overline{B}\cdot
\overline{C}
}
\]

যাচাই হলো।

% =========================================================
% Three-variable Second Theorem Truth Table
% =========================================================

\subsubsection*{তিন চলকের দ্বিতীয় উপপাদ্য যাচাই}

প্রমাণ করতে হবে—

\[
\overline{ABC}
=
\overline{A}+
\overline{B}+
\overline{C}
\]

\begin{center}
\scriptsize
\resizebox{\textwidth}{!}{%
\begin{tabular}{|c|c|c|c|c|c|}
\hline
\rowcolor{ICTLightBlue}
\(A\) &
\(B\) &
\(C\) &
\(ABC\) &
\(\overline{ABC}\) &
\(\overline{A}+\overline{B}+\overline{C}\) \\
\hline
0 & 0 & 0 & 0 & 1 & 1 \\
\hline
0 & 0 & 1 & 0 & 1 & 1 \\
\hline
0 & 1 & 0 & 0 & 1 & 1 \\
\hline
0 & 1 & 1 & 0 & 1 & 1 \\
\hline
1 & 0 & 0 & 0 & 1 & 1 \\
\hline
1 & 0 & 1 & 0 & 1 & 1 \\
\hline
1 & 1 & 0 & 0 & 1 & 1 \\
\hline
1 & 1 & 1 & 1 & 0 & 0 \\
\hline
\end{tabular}%
}
\end{center}

এখানেও শেষের দুটি কলামের মান প্রতিটি সারিতে সমান।

অতএব,

\[
\boxed{
\overline{ABC}
=
\overline{A}+
\overline{B}+
\overline{C}
}
\]

যাচাই হলো।

% =========================================================
% General n-variable Formula
% =========================================================

\subsubsection*{\(n\) সংখ্যক চলকের জন্য সাধারণ সূত্র}

ডি-মরগ্যানের উপপাদ্য যেকোনো সংখ্যক বুলিয়ান চলকের
ক্ষেত্রেই প্রযোজ্য।

\textbf{\textcolor{ICTOrange}{প্রথম উপপাদ্যের সাধারণ রূপ:}}

\[
\boxed{
\overline{
A_1+A_2+A_3+\cdots+A_n
}
=
\overline{A_1}\cdot
\overline{A_2}\cdot
\overline{A_3}\cdots
\overline{A_n}
}
\]

অর্থাৎ,
\textbf{\textenglish{OR}-এর পূরক =
প্রত্যেক চলকের পূরকের \textenglish{AND}}।

\textbf{\textcolor{ICTOrange}{দ্বিতীয় উপপাদ্যের সাধারণ রূপ:}}

\[
\boxed{
\overline{
A_1A_2A_3\cdots A_n
}
=
\overline{A_1}
+
\overline{A_2}
+
\overline{A_3}
+\cdots+
\overline{A_n}
}
\]

অর্থাৎ,
\textbf{\textenglish{AND}-এর পূরক =
প্রত্যেক চলকের পূরকের \textenglish{OR}}।

% =========================================================
% De Morgan Shortcut
% =========================================================

\par\begingroup

\textbf{ডি-মরগ্যানের সূত্র মনে রাখার সহজ কৌশল:}

কোনো রাশির বাইরের \textenglish{Complement}
ভেতরে প্রবেশ করলে—

\begin{enumerate}[leftmargin=*,itemsep=1pt]
    \item প্রতিটি চলকের ওপর \textenglish{Complement} বসবে।
    \item \textenglish{OR} \((+)\) পরিবর্তিত হয়ে
    \textenglish{AND} \((\cdot)\) হবে।
    \item \textenglish{AND} \((\cdot)\) পরিবর্তিত হয়ে
    \textenglish{OR} \((+)\) হবে।
\end{enumerate}

যেমন—

\[
\overline{A+B+C}
\;\Rightarrow\;
\overline{A}\cdot
\overline{B}\cdot
\overline{C}
\]

এবং,

\[
\overline{ABC}
\;\Rightarrow\;
\overline{A}+
\overline{B}+
\overline{C}
\]

\textbf{মনে রাখবে:}
ডি-মরগ্যানের উপপাদ্য শুধু দুইটি চলকের জন্য নয়;
দুই, তিন বা যেকোনো \(n\) সংখ্যক চলকের ক্ষেত্রেই প্রযোজ্য।
\par\endgroup


% =========================================================
% 3.8.10 Importance of Boolean Algebra
% =========================================================

\subsection{বুলিয়ান অ্যালজেবরার প্রয়োজনীয়তা}

ডিজিটাল সার্কিটে বিভিন্ন ইনপুটের ভিত্তিতে নির্দিষ্ট আউটপুট
পাওয়ার জন্য লজিক গেট ব্যবহৃত হয়।
এসব গেটের কার্যপ্রণালি বুলিয়ান অ্যালজেবরার সাহায্যে
প্রকাশ ও বিশ্লেষণ করা যায়।

বুলিয়ান রাশি সরল করলে—

\begin{enumerate}[leftmargin=*,itemsep=2pt]
    \item লজিক সার্কিটের গঠন সহজ হয়।
    \item প্রয়োজনীয় লজিক গেটের সংখ্যা কমানো যায়।
    \item সার্কিটের আকার ও বাস্তবায়ন ব্যয় কমতে পারে।
    \item সার্কিট বিশ্লেষণ, নকশা ও বাস্তবায়ন সহজ হয়।
\end{enumerate}

\par\begingroup

\textbf{এক নজরে:}

বুলিয়ান চলকের মান—

\[
A\in\{0,1\}
\]

মৌলিক ক্রিয়া—

\[
\boxed{
\text{\textenglish{OR,\ AND,\ NOT}}
}
\]

\(n\)টি ইনপুট চলক থাকলে সম্ভাব্য ইনপুট সমন্বয়ের সংখ্যা—

\[
\boxed{2^n}
\]

বুলিয়ান সূত্র ব্যবহার করে জটিল রাশি সরল করলে
প্রয়োজনীয় \textenglish{Logic Gate}-এর সংখ্যা কমানো যায়।
\par\endgroup

\subsection{দ্বৈত রূপ, পূরক ও সরলীকরণের যুক্তি}
দ্বৈত রূপ বের করতে অপারেশন ও ধ্রুবক বদলায়, চলকের ওপর নিজের থেকে পূরক বসে না। রাশির পূরক নির্ণয়ে সম্পূর্ণ রাশির মান ০ থেকে ১ বা ১ থেকে ০ হয়। এই দুই প্রক্রিয়া পৃথক।

\chThreeExample{১}{দ্বৈত রূপ এবং পূরক}
$F=A+BC$-এর দ্বৈত রূপ ও পূরক বের করো।
\chThreeAnswer দ্বৈত রূপ $A(B+C)$। ডি-মরগ্যান প্রয়োগে পূরক $\overline F=\overline A(\overline B+\overline C)$। দ্বৈত রূপ এবং মূল রাশির মান সব ইনপুটে একই হবে—এমন দাবি নেই; বৈধ \emph{সূত্রের} দ্বৈত সূত্রও বৈধ।

\smallhead{প্রয়োজনীয় সরলীকরণ সূত্রের কারণ}
\[\begin{aligned}
A+\overline AB&=(A+\overline A)(A+B)=A+B,\\
(A+B)(A+\overline B)&=A+B\overline B=A.
\end{aligned}\]
এখানে প্রথম লাইনে $A+BC=(A+B)(A+C)$, দ্বিতীয় লাইনে তার বিপরীত ব্যবহার হয়েছে। সাধারণ বীজগণিতের মতো $1+1=2$ ধরে হিসাব করলে ভুল হবে।

\smallhead{Consensus সূত্র}
\[AB+\overline AC+BC=AB+\overline AC.\]
কারণ $BC=BC(A+\overline A)=ABC+\overline ABC$। প্রথম নতুন পদ $AB$-এর এবং দ্বিতীয়টি $\overline AC$-এর মধ্যে আগেই অন্তর্ভুক্ত। ফলে $BC$ নতুন কোনো ইনপুটে আউটপুট ১ করছে না। দ্বৈত সূত্র $(A+B)(\overline A+C)(B+C)=(A+B)(\overline A+C)$।

\chThreeExample{২}{সরলীকরণে বাস্তব সিদ্ধান্ত অপরিবর্তিত}
এক যন্ত্র চালু হয় যদি প্রধান সুইচ $A$ চালু থাকে, অথবা $A$ ও অনুমতি সুইচ $B$ দুটোই চালু থাকে।
\chThreeAnswer $F=A+AB=A(1+B)=A$। দ্বিতীয় শর্ত প্রথমটির মধ্যেই অন্তর্ভুক্ত; $B$-এর জন্য আলাদা গেট দরকার নেই। সরলীকরণে শর্তের ফল বদলায়নি।

\chThreeExample{৩}{ডি-মরগ্যানে পূরকের সীমা}
$F=\overline{A+\overline{BC}}$ সরল করো।
\chThreeAnswer $F=\overline A\,\overline{\overline{BC}}=\overline ABC$। ভেতরের পুরো $BC$-এর দ্বৈত পূরক উঠে যায়; শুধু B-এর পূরক তোলা যাবে না।

\chThreeExample{৪}{ধারণাভিত্তিক MCQ}
বুলিয়ান বণ্টন সূত্র কোনটি? (ক) $A+BC=(A+B)(A+C)$, (খ) $A+A=2A$, (গ) $A+1=A$, (ঘ) $A\overline A=1$।
\chThreeAnswer (ক)। (খ), (গ), (ঘ)-তে যথাক্রমে $A$, $1$ ও $0$ সঠিক ফল।


\providecommand{\chLogicExample}[3]{\chThreeExample{#2}{#3}\label{ex:#1}}
\subsection{ধরন অনুযায়ী বুলিয়ান রাশির সমাধান}
প্রথমে রাশির পূরকের সীমা এবং বন্ধনী দেখো। তারপর কোন সূত্রে পদ কমবে তা নির্বাচন করো। নিচে প্রতিটি ধরনের দুইটি সমাধান রয়েছে; সমাধানের পরে সূত্রের কারণ বলা হয়েছে। বুলিয়ানে $+$ হলো OR, পাশাপাশি লেখা চলক বা $\cdot$ হলো AND।
\Needspace{12\baselineskip}
\subsubsection*{ধরন ১: অপারেশনের ক্রম ও মান বসানো}
আগে বন্ধনী ও পূরক, তারপর AND, সবশেষে OR করো। সাধারণ যোগের carry এখানে নেই।
\par\smallskip\noindent\begin{minipage}{\linewidth}
\chLogicExample{BA01.1}{১.১}{একই ১-এর দুই রকম যোগ}
$A=B=1,C=0$ হলে $F=A+B\overline C$ এবং বাইনারি যোগ $A+B$-এর ফল তুলনা করো.
\chThreeAnswer $\overline C=1$, তাই $B\overline C=1$ এবং বুলিয়ান $F=1+1=1$। কিন্তু বাইনারি গাণিতিক যোগে $1_2+1_2=10_2$। অপারেশন বোঝার আগে সংখ্যার চিহ্ন দেখে যোগ করা যাবে না।

\end{minipage}\par\smallskip
\par\smallskip\noindent\begin{minipage}{\linewidth}
\chLogicExample{BA01.2}{১.২}{আংশিক পূরকের মান}
$A=0,B=1,C=1$ হলে $F=(A+\overline B)C+\overline{AC}$ নির্ণয় করো.
\chThreeAnswer $\overline B=0$, $(A+\overline B)C=0$ এবং $AC=0$ থেকে $\overline{AC}=1$। তাই $F=0+1=1$। শেষ পূরকটি পুরো $AC$-এর।

\end{minipage}\par\smallskip
\Needspace{12\baselineskip}
\subsubsection*{ধরন ২: দ্বৈত রূপ নির্ণয়}
AND ও OR পরস্পর বদলাও; ০ ও ১ বদলাও। চলক এবং চলকের পূরক যেমন আছে তেমন থাকবে। দ্বৈত রূপকে মূল রাশির পূরক ভাববে না।
\par\smallskip\noindent\begin{minipage}{\linewidth}
\chLogicExample{BA02.1}{২.১}{বন্ধনীসহ দ্বৈত রূপ}
$F=(A+\overline B)(C+0)$-এর দ্বৈত রূপ লেখো.
\chThreeAnswer $F^d=A\overline B+C\cdot1=A\overline B+C$। $\overline B$ একই থাকে। মূল রাশিটি $C(A+\overline B)$, তাই দ্বৈত রূপের সঙ্গে সব input-এ সমান হওয়ার প্রয়োজন নেই।

\end{minipage}\par\smallskip
\par\smallskip\noindent\begin{minipage}{\linewidth}
\chLogicExample{BA02.2}{২.২}{ধ্রুবকসহ সূত্রের দ্বৈত}
$A\cdot1+AB=A$ সূত্রের দ্বৈত লিখে যাচাই করো.
\chThreeAnswer দ্বৈত হলো $(A+0)(A+B)=A$। বামপক্ষ $A(A+B)=A+AB=A$। একটি সত্য সূত্রের দ্বৈতও সত্য; চলকের নাম বদলানোর প্রয়োজন নেই।

\end{minipage}\par\smallskip
\Needspace{12\baselineskip}
\subsubsection*{ধরন ৩: সম্পূর্ণ রাশির পূরক}
ডি-মরগ্যানে OR-এর পূরক AND হয়, AND-এর পূরক OR হয়। প্রতিটি পদের পূরক নিতে হবে; দ্বৈত পূরক উঠে যায়।
\par\smallskip\noindent\begin{minipage}{\linewidth}
\chLogicExample{BA03.1}{৩.১}{যোগফলের পূরক}
$Y=\overline AB+BC$ হলে $\overline Y$ নির্ণয় করো.
\chThreeAnswer $\overline Y=\overline{\overline AB}\,\overline{BC}=(A+\overline B)(\overline B+\overline C)$। একই চলক বারবার থাকলেও প্রতিটি পদের পূরক আলাদাভাবে নেওয়া হয়েছে।

\end{minipage}\par\smallskip
\par\smallskip\noindent\begin{minipage}{\linewidth}
\chLogicExample{BA03.2}{৩.২}{গুণফলের পূরক}
$Y=(A+\overline C)(\overline B+D)$ হলে $\overline Y$ বের করো.
\chThreeAnswer $\overline Y=\overline{A+\overline C}+\overline{\overline B+D}=\overline AC+B\overline D$। এখানে বাইরের AND বদলে OR হয়েছে।

\end{minipage}\par\smallskip
\Needspace{12\baselineskip}
\subsubsection*{ধরন ৪: সাধারণ গুণক নিয়ে পূরক জোড়া করা}
একই গুণক বের করো। অবশিষ্ট অংশে $X+\overline X=1$ বা $X\overline X=0$ থাকলে তা ব্যবহার করো।
\par\smallskip\noindent\begin{minipage}{\linewidth}
\chLogicExample{BA04.1}{৪.১}{এক চলক বাদ পড়া}
$F=\overline AB\overline C+\overline ABC+AD$ সরল করো.
\chThreeAnswer প্রথম দুই পদের $\overline AB$ এক; C শুধু দুটি সম্ভাবনাই প্রকাশ করছে।
\[\begin{aligned}F&=\overline AB\overline C+\overline ABC+AD\\
&=\overline AB(\overline C+C)+AD\\
&=\overline AB+AD\end{aligned}\]

\end{minipage}\par\smallskip
\par\smallskip\noindent\begin{minipage}{\linewidth}
\chLogicExample{BA04.2}{৪.২}{একাধিক পূরক জোড়া}
$F=A\overline B\overline C+A\overline BC+AB\overline C+ABC$ সরল করো.
\chThreeAnswer সব পদের A এক। B ও C-এর সব সম্ভাবনা থাকায় তাদের মান output নির্ধারণ করে না।
\[\begin{aligned}F&=A\overline B\overline C+A\overline BC+AB\overline C+ABC\\
&=A\overline B(\overline C+C)+AB(\overline C+C)\\
&=A(\overline B+B)\\
&=A\end{aligned}\]

\end{minipage}\par\smallskip
\Needspace{12\baselineskip}
\subsubsection*{ধরন ৫: বড় শর্তের মধ্যে ছোট শর্ত শোষিত হওয়া}
$X+XY=X$ এবং $X(X+Y)=X$। X একটি যৌগিক পদও হতে পারে।
\par\smallskip\noindent\begin{minipage}{\linewidth}
\chLogicExample{BA05.1}{৫.১}{যৌগিক পদের absorption}
$F=AC+ABC+ACD$ সরল করো.
\chThreeAnswer $1+B+D=1$। তাই বাড়তি অনুমতি B,D ফল বদলায় না।
\[\begin{aligned}F&=AC+ABC+ACD\\
&=AC(1+B+D)\\
&=AC\end{aligned}\]

\end{minipage}\par\smallskip
\par\smallskip\noindent\begin{minipage}{\linewidth}
\chLogicExample{BA05.2}{৫.২}{POS-এ absorption}
$F=(A+D)(A+D+B)(A+D+C)$ সরল করো.
\chThreeAnswer $X=A+D$ ধরলে $X(X+B)(X+C)=X$।
\[\begin{aligned}F&=(A+D)(A+D+B)(A+D+C)\\
&=(A+D)\end{aligned}\]

\end{minipage}\par\smallskip
\Needspace{12\baselineskip}
\subsubsection*{ধরন ৬: পূরকসহ elimination}
$X+\overline XY=X+Y$ এবং $X(\overline X+Y)=XY$। সরাসরি absorption-এর সঙ্গে এই সূত্রের পার্থক্য আছে।
\par\smallskip\noindent\begin{minipage}{\linewidth}
\chLogicExample{BA06.1}{৬.১}{যৌগিক Y নিয়ে elimination}
$F=\overline D+DAB$ সরল করো.
\chThreeAnswer $X=\overline D$, $Y=AB$। বণ্টন সূত্রে প্রথম গুণক ১ হয়।
\[\begin{aligned}F&=\overline D+DAB\\
&=(\overline D+D)(\overline D+AB)\\
&=\overline D+AB\end{aligned}\]

\end{minipage}\par\smallskip
\par\smallskip\noindent\begin{minipage}{\linewidth}
\chLogicExample{BA06.2}{৬.২}{AND শর্তে বিপরীত চলক}
$F=(A+C)(\overline A\,\overline C+B)$ সরল করো.
\chThreeAnswer $\overline{A+C}=\overline A\,\overline C$। X ও তার পূরকের গুণফল শূন্য।
\[\begin{aligned}F&=(A+C)(\overline A\,\overline C+B)\\
&=(A+C)\overline A\,\overline C+(A+C)B\\
&=B(A+C)\end{aligned}\]

\end{minipage}\par\smallskip
\Needspace{12\baselineskip}
\subsubsection*{ধরন ৭: POS প্রসারণ ও SOP-তে factor করা}
$X+YZ=(X+Y)(X+Z)$ এবং $(X+Y)(X+Z)=X+YZ$। সাধারণ বীজগণিতে প্রথম সূত্রটি একইভাবে চলে না।
\par\smallskip\noindent\begin{minipage}{\linewidth}
\chLogicExample{BA07.1}{৭.১}{সাধারণ OR পদ}
$F=(\overline A+B)(\overline A+C)$-কে SOP করো.
\chThreeAnswer $\overline A\,\overline A=\overline A$; তার সঙ্গে থাকা B,C পদ absorption-এ বাদ পড়ে।
\[\begin{aligned}F&=(\overline A+B)(\overline A+C)\\
&=\overline A+\overline AC+\overline AB+BC\\
&=\overline A+BC\end{aligned}\]

\end{minipage}\par\smallskip
\par\smallskip\noindent\begin{minipage}{\linewidth}
\chLogicExample{BA07.2}{৭.২}{একাধিক গুণকের factor}
$F=AB+CD$-কে POS-এ লেখো.
\chThreeAnswer দুইবার বণ্টন সূত্র প্রয়োগ করো। এই রূপের সব গুণক রাখতে হবে।
\[\begin{aligned}F&=AB+CD\\
&=(AB+C)(AB+D)\\
&=(A+C)(B+C)(A+D)(B+D)\end{aligned}\]

\end{minipage}\par\smallskip
\Needspace{12\baselineskip}
\subsubsection*{ধরন ৮: একাধিক স্তরের পূরক}
সবচেয়ে বাইরের পূরক খুলে এক স্তর করে এগোও। প্রতিটি overbar কোন পদের ওপর আছে তা চিহ্নিত করো।
\par\smallskip\noindent\begin{minipage}{\linewidth}
\chLogicExample{BA08.1}{৮.১}{ভেতরের পূরক উঠে যাওয়া}
$F=\overline{\overline{AB}+C}$ সরল করো.
\chThreeAnswer ভেতরের AB দ্বৈত পূরক হয়ে ফিরে আসে।
\[\begin{aligned}F&=\overline{\overline{AB}+C}\\
&=\overline{\overline{AB}}\,\overline C\\
&=AB\overline C\end{aligned}\]

\end{minipage}\par\smallskip
\par\smallskip\noindent\begin{minipage}{\linewidth}
\chLogicExample{BA08.2}{৮.২}{বন্ধনীর ওপর দ্বিতীয় পূরক}
$F=\overline{(A+\overline B)\overline{C+D}}$ সরল করো.
\chThreeAnswer বাইরের AND-এর পূরক OR; দ্বিতীয় পদের দ্বৈত পূরক উঠে যায়।
\[\begin{aligned}F&=\overline{(A+\overline B)\overline{C+D}}\\
&=\overline{A+\overline B}+\overline{\overline{C+D}}\\
&=\overline AB+C+D\end{aligned}\]

\end{minipage}\par\smallskip
\Needspace{12\baselineskip}
\subsubsection*{ধরন ৯: ধ্রুবক ০ বা ১ পাওয়া}
পরস্পরবিরোধী প্রয়োজন একসঙ্গে AND হলে output ০। সব সম্ভাবনাকে OR করলে output ১।
\par\smallskip\noindent\begin{minipage}{\linewidth}
\chLogicExample{BA09.1}{৯.১}{দুটি বিপরীত শর্ত}
$F=(A+C)\overline A\,\overline C$ সরল করো.
\chThreeAnswer একই সঙ্গে A অথবা C-কে ১ এবং দুটিকেই ০ চাওয়া হয়েছে।
\[\begin{aligned}F&=(A+C)\overline A\,\overline C\\
&=A\overline A\,\overline C+C\overline A\,\overline C\\
&=0\end{aligned}\]

\end{minipage}\par\smallskip
\par\smallskip\noindent\begin{minipage}{\linewidth}
\chLogicExample{BA09.2}{৯.২}{সব সম্ভাবনায় output ১}
$F=AB+A\overline B+\overline AC+\overline A\,\overline C$ সরল করো.
\chThreeAnswer A=১ ও A=০—দুটির জন্যই অন্য চলকের সব সম্ভাবনা রাখা হয়েছে।
\[\begin{aligned}F&=AB+A\overline B+\overline AC+\overline A\,\overline C\\
&=A(B+\overline B)+\overline A(C+\overline C)\\
&=A+\overline A\\
&=1\end{aligned}\]

\end{minipage}\par\smallskip
\Needspace{12\baselineskip}
\subsubsection*{ধরন ১০: Consensus-এ অপ্রয়োজনীয় পদ বাদ দেওয়া}
এক চলক বিপরীত অবস্থায় দুই পদে থাকলে তাদের অন্য অংশের গুণফল অনেক সময় আগেই অন্তর্ভুক্ত থাকে। পদ বাদ দেওয়ার আগে সূত্রের তিন অংশ মেলাও।
\par\smallskip\noindent\begin{minipage}{\linewidth}
\chLogicExample{BA10.1}{১০.১}{SOP consensus}
$F=AC+\overline A D+CD$ সরল করো.
\chThreeAnswer $CD=CD(A+\overline A)=ACD+\overline ACD$; নতুন পদদুটি প্রথম দুই পদের মধ্যে আছে।
\[\begin{aligned}F&=AC+\overline A D+CD\\
&=AC+\overline A D+ACD+\overline ACD\\
&=AC+\overline A D\end{aligned}\]

\end{minipage}\par\smallskip
\par\smallskip\noindent\begin{minipage}{\linewidth}
\chLogicExample{BA10.2}{১০.২}{POS consensus}
$F=(A+C)(\overline A+D)(C+D)$ সরল করো.
\chThreeAnswer প্রথম দুই বন্ধনী ১ হলে তৃতীয়টিও ১: A=০ হলে C=১ লাগবে; A=১ হলে D=১ লাগবে।
\[\begin{aligned}F&=(A+C)(\overline A+D)(C+D)\\
&=(A+C)(\overline A+D)\end{aligned}\]

\end{minipage}\par\smallskip
\Needspace{12\baselineskip}
\subsubsection*{ধরন ১১: রাশিতে XOR বা XNOR চেনা}
দুই চলকের বিপরীত মানের পদদুটি XOR; একই মানের পদদুটি XNOR। Gate-এর bubble থাকলে শেষ পূরকটিও দেখো।
\par\smallskip\noindent\begin{minipage}{\linewidth}
\chLogicExample{BA11.1}{১১.১}{POS থেকে XOR}
$F=(A+D)(\overline A+\overline D)$ সরল করে gate চিহ্নিত করো.
\chThreeAnswer $A\overline D+\overline AD=A\oplus D$।
\[\begin{aligned}F&=(A+D)(\overline A+\overline D)\\
&=A\overline A+A\overline D+D\overline A+D\overline D\\
&=A\overline D+\overline AD\end{aligned}\]

\end{minipage}\par\smallskip
\par\smallskip\noindent\begin{minipage}{\linewidth}
\chLogicExample{BA11.2}{১১.২}{XOR-এর পূরকে সমতা}
$F=\overline{B\overline C+\overline BC}$ সরল করো.
\chThreeAnswer এটি $\overline{B\oplus C}$ বা XNOR।
\[\begin{aligned}F&=\overline{B\overline C+\overline BC}\\
&=(\overline B+C)(B+\overline C)\\
&=\overline B\,\overline C+BC\end{aligned}\]

\end{minipage}\par\smallskip
\Needspace{12\baselineskip}
\subsubsection*{ধরন ১২: সূত্র প্রমাণ ও ভুল সমতা খণ্ডন}
সমতুল্যতা প্রমাণে সব input-এ output মিলতে হবে। একটি বিপরীত ফল পেলেই দাবিটি ভুল প্রমাণিত হয়।
\par\smallskip\noindent\begin{minipage}{\linewidth}
\chLogicExample{BA12.1}{১২.১}{সত্যক সারণিতে যৌগিক সূত্র}
$AB+A\overline B=A$ সত্যক সারণিতে দেখাও.
\chThreeAnswer প্রতিটি সারিতে শেষ দুই column সমান।
\begin{center}\begin{tabular}{|c|c|c|c|c|c|}\hline
$A$ & $B$ & $AB$ & $A\overline B$ & $AB+A\overline B$ & $A$\\\hline
0 & 0 & 0 & 0 & 0 & 0\\\hline
0 & 1 & 0 & 0 & 0 & 0\\\hline
1 & 0 & 0 & 1 & 1 & 1\\\hline
1 & 1 & 1 & 0 & 1 & 1\\\hline
\end{tabular}\end{center}

\end{minipage}\par\smallskip
\par\smallskip\noindent\begin{minipage}{\linewidth}
\chLogicExample{BA12.2}{১২.২}{একটি input দিয়ে ভুল ধরা}
এক শিক্ষার্থী লিখেছে $\overline{A+B}=\overline A+\overline B$। দাবিটি পরীক্ষা করো.
\chThreeAnswer $A=0,B=1$-এ বামপক্ষ ০, ডানপক্ষ $1+0=1$। তাই সমতা ভুল। সঠিক সূত্র $\overline{A+B}=\overline A\,\overline B$। শুধু 00 input পরীক্ষা করলে ভুলটি ধরা পড়ত না।

\end{minipage}\par\smallskip
\smallhead{ক--খ প্রশ্নের ধারণা}
বুলিয়ান $1+1=1$ কারণ অন্তত একটি সত্য শর্তে OR সত্য হয়। ডি-মরগ্যানের সূত্র দিয়ে বড় পূরক ভেঙে পৃথক input-এর পূরক এবং AND/OR দিয়ে সমতুল্য রাশি পাওয়া যায়। সরলীকরণে সত্যক সারণি অপরিবর্তিত রেখে প্রয়োজনীয় gate ও সংযোগ কমতে পারে। দ্বৈত রূপ এবং পূরক পৃথক প্রক্রিয়া।

